\documentclass[12pt, a4paper]{article}

% --- تنظیمات کلی و حاشیه‌ها ---
\usepackage[top=2.5cm, bottom=2.5cm, left=2cm, right=2cm]{geometry}
\usepackage{fontspec}

% --- تنظیمات زبان و فونت با Babel ---
\usepackage[persian, bidi=basic, provide=*]{babel}

\babelprovide[import, onchar=ids fonts]{persian}
\babelprovide[import, onchar=ids fonts]{english}

% تنظیم فونت اصلی روی امیری و انگلیسی روی Noto Sans
\babelfont{rm}{Amiri}
\babelfont[persian]{rm}{Amiri}
\babelfont[english]{rm}{Noto Sans}

\usepackage{enumitem}
\setlist[itemize]{label=-}

\usepackage{tabularx}
\usepackage{xcolor}

% تعریف یک دستور ساده برای نوشتن کلمات انگلیسی در متن فارسی
\newcommand{\lr}[1]{\foreignlanguage{english}{#1}}

% --- تنظیمات باکس‌های کد و خروجی ---
\usepackage{tcolorbox}
\tcbuselibrary{listings}

% محیط برای کدهای برنامه نویسی (طراحی روشن، شبیه ادیتور کد)
\newtcblisting{codebox}{
  colback=gray!10!white, colupper=black, colframe=blue!40!black,
  listing only, boxrule=0.8pt, arc=4pt, left=2mm,
  listing options={
    basicstyle=\ttfamily\small\color{black}, 
    columns=fullflexible, 
    breaklines=true, 
     breakatwhitespace=true,
    postbreak=\mbox{\textcolor{red}{$\hookrightarrow$}\space}
  }
}

% محیط برای خروجی‌های ترمینال (طراحی تیره، شبیه ترمینال اوبونتو)
\newtcblisting{terminalbox}{
  colback=black!85, colupper=green!80!white, colframe=black!70,
  listing only, boxrule=0.5pt, arc=4pt, left=2mm,
  listing options={
    basicstyle=\ttfamily\footnotesize\color{green!80!white}, 
    columns=fullflexible, 
    breaklines=true,
    breakatwhitespace=true,
    postbreak=\mbox{\textcolor{yellow}{$\hookrightarrow$}\space}
  }
}

% بسته hyperref برای لینک‌دار شدن فهرست مطالب (باید در انتهای preamble باشد)
% تغییر داده شد تا لینک‌های فهرست به صورت رنگی و کاملاً قابل کلیک باشند
\usepackage[colorlinks=true, linkcolor=blue, urlcolor=blue, citecolor=blue]{hyperref}

\begin{document}
\thispagestyle{empty}

% --- شناسنامه گزارش ---
\begin{center}
\renewcommand{\arraystretch}{1.8}
\begin{tabularx}{\textwidth}{|c|c|c|X|c|c|}
\hline
\multicolumn{6}{|c|}{\textbf{شناسنامه‌ی گزارش}} \\
\hline
\multicolumn{2}{|c|}{\textbf{عنوان گزارش}} & \multicolumn{4}{l|}{راه‌اندازی اولیه و پایش \lr{RSSI} در پروژه \lr{RuView}} \\
\hline
\multicolumn{6}{|r|}{\textbf{چکیده}} \\
\hline
\multicolumn{6}{|p{\dimexpr\textwidth-2\tabcolsep-2\arrayrulewidth\relax}|}{
در این گزارش، مراحل استقرار پروژه \lr{RuView} برای سنجش وای‌فای مستند شده است. تمرکز اصلی بر استخراج اطلاعات \lr{RSSI} با استفاده از ماژول \lr{v1} و پیاده‌سازی سیستم طبقه‌بندی حضور در محیط لینوکس است.
} \\
\hline
\textbf{نسخه} & \textbf{تاریخ} & \textbf{نگارش کننده} & \multicolumn{1}{c|}{\textbf{تغییرات}} & \textbf{تایید} & \textbf{مدیر} \\
\hline
1 & 1405/05/06 & امیرعلی یزدیانی & $\bullet$ تدوین مستند و راه‌اندازی & & \\
\hline
\end{tabularx}
\end{center}

\vspace{1cm}
% --- فهرست مطالب خودکار و لینک‌دار ---
\tableofcontents
\newpage

\section{چشم‌انداز و معرفی پروژه}
پروژه \lr{RuView} یک پلتفرم پیشرفته برای تبدیل سیگنال‌های عادی وای‌فای به سیستم سنجش فضایی است. هدف ما این است که بدون استفاده از دوربین یا ابزارهای پوشیدنی، حضور افراد، حرکت و حتی علائم حیاتی را تنها از طریق تحلیل اختلالات امواج رادیویی تشخیص دهیم. در فاز فعلی، هدفِ ما یادگیری روند «جمع‌آوری داده»، «استخراج ویژگی» و «طبقه‌بندی وضعیت حضور» است.

\section{آماده‌سازی محیط}
برای پیاده‌سازی بخش \lr{v1}، اول از همه تمامی بخش‌های پروژه را از \lr{repository} گیت‌هاب برمی‌داریم:

\begin{otherlanguage}{english}
\begin{terminalbox}
git clone https://github.com/ruvnet/RuView.git
cd RuView
\end{terminalbox}
\end{otherlanguage}

بعد از آن باید یک محیط مجازی پایتون درست کنیم و پیش‌نیازهایی که برایمان گذاشته شده را به علاوه برخی پیش‌نیازهای دیگر نصب کنیم:

\begin{otherlanguage}{english}
\begin{terminalbox}
cd archive/v1
python3 -m venv .venv
source .venv/bin/activate
.venv/bin/pip install -r requirements-lock.txt
.venv/bin/pip install psutil fastapi uvicorn
\end{terminalbox}
\end{otherlanguage}

با توجه به این که کدهای اولیه برای \lr{Windows} نوشته شده‌اند، برخی تغییرات باید انجام شوند تا برای محیط لینوکس آماده‌سازی شوند. این دو تغییر را در آن اعمال و فایل را ذخیره کنید:
\begin{itemize}
    \item وارد کردن ماژول را تغییر دهید: \lr{WindowsWifiCollector} را پاک کرده و به جای آن \lr{LinuxWifiCollector} را بنویسید.
    \item کد اجرای کالکتور را پیدا کرده و اینترفیس خودتان را جایگزین کنید:
\end{itemize}

\begin{otherlanguage}{english}
\begin{codebox}
collector = LinuxWifiCollector(interface="wlx002e2d2091fa", sample_rate_hz=2.0)
\end{codebox}
\end{otherlanguage}

\section{استخراج اطلاعات \lr{RSSI}}
برای خواندن \lr{CSI} به محاسبات زیادی و سخت‌افزار مربوطه (\lr{ESP32}) نیاز داریم، ولی می‌توانیم از \lr{RSSI} استفاده کنیم تا با استفاده از \lr{V1} مقداری از اطلاعات کانال را به دست آوریم و پردازش‌های اولیه‌ای انجام دهیم. 

ابتدا با استفاده از دستور \lr{iwconfig} نام اینترفیسی که با آن به شبکه وصل هستیم را پیدا می‌کنیم:

\begin{otherlanguage}{english}
\begin{terminalbox}
epo@hpds:~/Desktop/Amirali/WifiSensing/ruview/archive/v1$ iwconfig
lo        no wireless extensions.

eno1      no wireless extensions.

kvnet     no wireless extensions.

wlx002e2d2091fa  IEEE 802.11  ESSID:"Sharif-WiFi"  
          Mode:Managed  Frequency:2.412 GHz  Access Point: 04:18:D6:E2:8D:35   
          Bit Rate=6.5 Mb/s   Tx-Power=20 dBm   
          Retry short limit:7   RTS thr=2347 B   Fragment thr:off
          Power Management:off
          Link Quality=54/70  Signal level=-56 dBm  
          Rx invalid nwid:0  Rx invalid crypt:0  Rx invalid frag:0
          Tx excessive retries:0  Invalid misc:7399   Missed beacon:0
\end{terminalbox}
\end{otherlanguage}

سپس با استفاده از کد زیر پارامترهای \lr{RSSI} را ثبت می‌کنیم:

\begin{otherlanguage}{english}
\begin{codebox}
### Section 1
import sys; sys.path.insert(0, 'archive/v1')
from src.sensing.rssi_collector import LinuxWifiCollector

c = LinuxWifiCollector(interface='wlx002e2d2091fa')
s = c.collect_once()
print(f'RSSI: {s.rssi_dbm} dBm, Quality: {s.link_quality:.0%}')
\end{codebox}
\end{otherlanguage}

خروجی آن نشان می‌دهد که توانستیم با استفاده از \lr{WiFi} کامپیوتر اطلاعاتی از سیگنال در اطراف خود پیدا کنیم:

\begin{otherlanguage}{english}
\begin{terminalbox}
epo@hpds:~/Desktop/Amirali/WifiSensing/ruview/archive/v1$ PYTHONPATH=.. .venv/bin/python test36.py
RSSI: -52.0 dBm, Quality: 83%
\end{terminalbox}
\end{otherlanguage}

\section{پردازش داده‌ها}
حال با استفاده از اطلاعاتی که به دست آورده‌ایم، به کمک \lr{RssiFeatureExtractor} و \lr{PresenceClassifier} اطلاعات را پردازش کرده و خروجی‌های خواسته شده را به دست می‌آوریم. 

این روند و خروجی آن در کد زیر مشخص است:

\begin{otherlanguage}{english}
\begin{codebox}
import sys, time
from src.sensing.rssi_collector import LinuxWifiCollector
from src.sensing.feature_extractor import RssiFeatureExtractor
from src.sensing.classifier import PresenceClassifier

collector = LinuxWifiCollector(interface='wlx002e2d2091fa', sample_rate_hz=2.0)
extractor = RssiFeatureExtractor(window_seconds=15.0)
classifier = PresenceClassifier(presence_variance_threshold=0.3)

collector.start()
print('Collecting 15 seconds of RSSI data... (Walk around to test motion!)')
time.sleep(15)
collector.stop()

samples = collector.get_samples()
if not samples:
    print('No samples collected!')
else:
    features = extractor.extract(samples)
    result = classifier.classify(features)
    
    print(f'Samples: {len(samples)}')
    print(f'RSSI mean: {features.mean:.1f} dBm')
    print(f'Variance: {features.variance:.4f}')
    print(f'Motion Energy: {features.motion_band_power:.4f}')
    print(f'Verdict: {getattr(result, "motion_level", result)}')
\end{codebox}
\end{otherlanguage}

خروجی اجرای کد به شکل زیر است:

\begin{otherlanguage}{english}
\begin{terminalbox}
epo@hpds:~/Desktop/Amirali/WifiSensing/ruview/archive/v1$ PYTHONPATH=.. .venv/bin/python test36.py
Collecting 15 seconds of RSSI data... (Walk around to test motion!)
Samples: 30
RSSI mean: -44.3 dBm
Variance: 13.1678
Motion Energy: 14.8927
Verdict: MotionLevel.ACTIVE
\end{terminalbox}
\end{otherlanguage}

با استفاده از \lr{RSSI} به دست آمده در ۱۵ ثانیه، توانستیم واریانس، انرژی حرکت و یک حدس اولیه برای حرکت کردن یا نکردن به دست بیاوریم. 

نکته‌ای که وجود دارد این است که حدس حرکت کردن صرفاً با یک آستانه (\lr{threshold}) محاسبه می‌شود که معیار خیلی دقیقی نیست، ولی در این مرحله بیشتر از این کاری از دست ما برنمی‌آید.

\section{پایش لحظه‌ای (\lr{Live Monitoring})}
همچنین با نوشتن اسکریپت \lr{live\_monitoring\_for\_linux} می‌توانیم در تمامی لحظات اطلاعاتی که کامپیوتر از \lr{WiFi} دریافت می‌کند را بررسی کنیم. 

در این بخش با تغییر دادن مقدار آستانه حرکت (\lr{motion threshold})، به یک معیار بسیار بهتر برای تشخیص حرکت دست پیدا می‌کنیم و می‌توانیم اطلاعات را به صورت پیوسته ثبت و اندازه‌گیری کنیم:

\begin{otherlanguage}{english}
\begin{codebox}
import time
from src.sensing.rssi_collector import LinuxWifiCollector
from src.sensing.feature_extractor import RssiFeatureExtractor
from src.sensing.classifier import PresenceClassifier

collector = LinuxWifiCollector(interface='wlx002e2d2091fa', sample_rate_hz=2.0)
extractor = RssiFeatureExtractor(window_seconds=15.0)
classifier = PresenceClassifier(presence_variance_threshold=0.3)

collector.start()
print("Starting Live Monitor on wlx002e2d2091fa... (Walk around to test!)")
print("Press Ctrl+C to stop.")
print("-" * 60)

try:
    while True:
        time.sleep(3)
        samples = collector.get_samples()
        
        if len(samples) > 5:
            features = extractor.extract(samples)
            result = classifier.classify(features)
            
            # Format the time and result to match the tutorial's output
            current_time = time.strftime("%H:%M:%S")
            verdict = getattr(result, "motion_level", result)
            
            print(f"[{current_time}] RSSI= {features.mean:.1f}dBm  motion={features.motion_band_power:.4f} var={features.variance:.3f} => {verdict}")
        else:
            print("Waiting for enough samples...")
            
except KeyboardInterrupt:
    print("\nStopping Live Monitor...")
finally:
    collector.stop()
\end{codebox}
\end{otherlanguage}

خروجی اجرای این اسکریپت در ترمینال به صورت زیر خواهد بود که وضعیت حضور و حرکت را در لحظه نشان می‌دهد:

\begin{otherlanguage}{english}
\begin{terminalbox}
epo@hpds:~/Desktop/Amirali/WifiSensing/ruview/archive/v1$ PYTHONPATH=.. .venv/bin/python live_monitoring_for_linux.py 
Starting Live Monitor on wlx002e2d2091fa... (Walk around to test!)
Press Ctrl+C to stop.
------------------------------------------------------------
[14:51:27] RSSI= -52.0dBm  motion=0.0000 var=0.000 => MotionLevel.ABSENT
[14:51:30] RSSI= -52.0dBm  motion=0.0000 var=0.000 => MotionLevel.ABSENT
[14:51:33] RSSI= -62.9dBm  motion=13.6989 var=2322.173 => MotionLevel.ACTIVE
[14:51:36] RSSI= -57.9dBm  motion=5641.5963 var=1717.193 => MotionLevel.ACTIVE
[14:51:39] RSSI= -55.2dBm  motion=11071.2144 var=1455.151 => MotionLevel.ACTIVE
[14:51:42] RSSI= -54.0dBm  motion=5225.5977 var=1471.379 => MotionLevel.ACTIVE
[14:51:45] RSSI= -52.2dBm  motion=133.8186 var=1494.626 => MotionLevel.ACTIVE
[14:51:48] RSSI= -43.8dBm  motion=18.0103 var=7.340 => MotionLevel.ACTIVE
[14:51:51] RSSI= -43.4dBm  motion=7.0809 var=4.861 => MotionLevel.PRESENT_STILL
[14:51:54] RSSI= -43.0dBm  motion=4.1193 var=4.447 => MotionLevel.PRESENT_STILL
[14:51:57] RSSI= -42.7dBm  motion=3.9508 var=2.023 => MotionLevel.PRESENT_STILL
[14:52:00] RSSI= -44.2dBm  motion=3.7641 var=10.855 => MotionLevel.PRESENT_STILL
[14:52:03] RSSI= -45.8dBm  motion=11.0677 var=15.289 => MotionLevel.ACTIVE
[14:52:06] RSSI= -47.3dBm  motion=8.1673 var=15.390 => MotionLevel.PRESENT_STILL
[14:52:09] RSSI= -48.8dBm  motion=7.1203 var=8.668 => MotionLevel.PRESENT_STILL
[14:52:12] RSSI= -50.0dBm  motion=11.9911 var=2.414 => MotionLevel.ACTIVE
[14:52:15] RSSI= -57.0dBm  motion=7.9907 var=1415.034 => MotionLevel.PRESENT_STILL
[14:52:18] RSSI= -57.1dBm  motion=2486.7691 var=1413.955 => MotionLevel.ACTIVE
[14:52:21] RSSI= -57.2dBm  motion=10558.1736 var=1411.978 => MotionLevel.ACTIVE
[14:52:24] RSSI= -57.1dBm  motion=6777.4354 var=1414.507 => MotionLevel.ACTIVE
[14:52:27] RSSI= -56.9dBm  motion=347.3317 var=1417.844 => MotionLevel.ACTIVE
\end{terminalbox}
\end{otherlanguage}

\subsection{نتیجه‌گیری فاز اول و پایان کار \lr{V1}}
حال بدون استفاده از سخت‌افزار خارجی توانستیم با استفاده از \lr{RSSI} به موارد خواسته شده برسیم. با استفاده از \lr{RSSI} به ضربان قلب و تنفس و تعداد افراد حاضر نیز می‌توان رسید، ولی برای این کار نیاز به استفاده از \lr{Rust} (که زبان سریع‌تری است) داریم که در بخش دو به آن می‌پردازیم. 

با اتمام این بخش به اتمام \lr{v1} می‌رسیم که همان‌طور که نویسنده ریپازیتوری گفته، به انتهای خود رسیده و ادامه ندارد.

\newpage
\section{بخش دوم: توسعه پیشرفته با زبان \lr{Rust}}

\subsection{دلیل مهاجرت به زبان \lr{Rust}}
ابتدا دلیل این که به این بخش می‌آییم را بررسی می‌کنیم. اطلاعاتی که از \lr{WiFi} می‌گیریم زیاد و پردازش نشده هستند؛ همچنین در یک برنامه بلادرنگ (\lr{Real-time Application}) که داریم استفاده می‌کنیم، نیاز به سرعت بالایی داریم. این‌ها مواردی هستند که در زبان \lr{Python} در اختیار نداشتیم؛ برای همین به سراغ زبان \lr{Rust} می‌آییم.

\subsection{نصب و راه‌اندازی محیط \lr{Rust}}
ابتدا به پیاده‌سازی و نصب \lr{Rust} می‌پردازیم که کار زمان‌بری است؛ ولی وقتی که دفعه اول انجام شود، دیگر در سری‌های بعدی با زمان بسیار کمتری انجام می‌شود و همچنین سرعت بسیار خوبی در اجرا دارد.

ابتدا از پوشه \lr{v1} خارج شدیم و وارد پوشه نسخه \lr{Rust} شدیم:

\begin{otherlanguage}{english}
\begin{terminalbox}
cd ../v2
\end{terminalbox}
\end{otherlanguage}

وقتی می‌خواستیم کدها را بخوانیم یا بیلد بگیریم، متوجه شدیم بعضی از پوشه‌های حیاتی (مثل کتابخانه‌های جانبی که نویسنده لینک کرده بود) کاملاً خالی هستند.

\textbf{راه‌حل:} دلیل این اتفاق این بود که گیت هنگام کلون کردنِ اولیه، ساب‌ماژول‌ها (\lr{Submodules}) را دانلود نکرده بود. با دستور زیر به گیت گفتیم تمام وابستگی‌های لینک‌شده را نیز دانلود کند:

\begin{otherlanguage}{english}
\begin{terminalbox}
git submodule update --init --recursive
\end{terminalbox}
\end{otherlanguage}

\subsection{تلاش اول برای کامپایل و خطای \lr{OpenBLAS}}
بعد از دریافت ساب‌ماژول‌ها، با ابزار \lr{Cargo} (پکیج‌منیجر \lr{Rust}) دستور کامپایل را اجرا کردیم:

\begin{otherlanguage}{english}
\begin{terminalbox}
cargo build --release -p wifi-densepose-sensing-server
\end{terminalbox}
\end{otherlanguage}

\textbf{راه‌حل:} این خطا به خاطر این بود که سرور \lr{Rust} برای محاسبات سنگین ریاضی و ماتریسی روی سیگنال‌ها، به کتابخانه \lr{C} به نام \lr{OpenBLAS} نیاز داشت که روی سیستم لینوکس شما نصب نبود. با نصب آن، مشکل کاملاً برطرف شد:

\begin{otherlanguage}{english}
\begin{terminalbox}
sudo apt-get update
sudo apt-get install libopenblas-dev pkg-config
\end{terminalbox}
\end{otherlanguage}

\subsection{بیلد نهایی و اجرای سرور}
بعد از نصب پیش‌نیازهای سیستمی، دوباره دستور بیلد را اجرا کردیم. این بار \lr{Cargo} تمام \lr{Crate}ها (پکیج‌های \lr{Rust}) را بدون مشکل دانلود و کامپایل کرد. (چون از فلگ \lr{--release} استفاده کردیم، کدها بهینه‌سازی شدند تا بالاترین سرعت ممکن را در اجرا داشته باشند).

حالا فایل اجرایی سرور ساخته شده و در مسیر \lr{target/release/} قرار داشت.

\textbf{چالش:} به صورت پیش‌فرض، اگر فقط سرور را اجرا می‌کردیم، سیستم به دنبال یک برد \lr{ESP32} می‌گشت تا دیتای \lr{CSI} را دریافت کند.

\textbf{راه‌حل:} برای اینکه به سرور بفهمانیم برد نداریم و می‌خواهیم دیتای \lr{RSSI} را از طریق دانگل وای‌فای لینوکسی خودمان بخواند، از آرگومان‌های \lr{--source} و \lr{--interface} استفاده کردیم و نام دقیق کارت شبکه شما را به آن پاس دادیم:

\begin{otherlanguage}{english}
\begin{terminalbox}
./target/release/sensing-server --source linux --interface wlx002e2d2091fa
\end{terminalbox}
\end{otherlanguage}

\subsection{اجرای سرور با \lr{API} اصلی و بررسی علائم حیاتی}
علاوه بر حالت قبل، با استفاده از دستور زیر می‌توانیم سرور را با استفاده از داده‌های \lr{RSSI} و روی حالت خودکار (\lr{auto}) اجرا کنیم تا سیستم به طور کامل راه‌اندازی شود:

\begin{otherlanguage}{english}
\begin{terminalbox}
./target/release/sensing-server --source auto
\end{terminalbox}
\end{otherlanguage}

پس از اجرای این دستور، در بخش نتایج به خوبی مشاهده می‌کنیم که سیستم چگونه حضور افراد را با درصدهای تشابه و اطمینان مختلف تشخیص می‌دهد. 

با این حال، برای استخراج اطلاعات دقیق‌تر مانند «تشخیص ضربان قلب» و «تنفس»، نیاز است که \lr{Access Point} (مودم یا روتر) و دانگل کامپیوتر در یک مسیر مستقیم و بدون مانع فیزیکی (\lr{Line of Sight}) نسبت به یکدیگر قرار داشته باشند. از آنجایی که در محیط آزمایش فعلی ما این شرایط به صورت ایده‌آل ممکن نیست، سیگنال‌ها با نویز همراه بوده و همین امر باعث می‌شود که در زمینه علائم حیاتی اطلاعات پرت و نامعتبری دریافت کنیم.

برای پایش و مشاهده‌ی خروجیِ علائم حیاتی از طریق \lr{API} لوکال، می‌توانیم از دستور \lr{curl} زیر در یک ترمینال جدید استفاده کنیم:

\begin{otherlanguage}{english}
\begin{terminalbox}
curl http://localhost:8080/api/v1/vital-signs
\end{terminalbox}
\end{otherlanguage}

\section{ارزیابی عملکرد و پردازش داده‌های از پیش‌ضبط‌شده (\lr{CSI}) در نسخه \lr{Rust}}
پس از آنکه در مرحله اجرای زنده (\lr{Live Sensing}) در لینوکس با خطای وابستگی به \lr{netsh} مواجه شدیم و با توجه به عدم دسترسی به برد سخت‌افزاری \lr{ESP32}، تصمیم گرفتیم به جای رها کردن معماری پیشرفته‌ی \lr{v2}، هسته‌ی پردازشی آن را از طریق داده‌های از پیش‌ضبط‌شده (آفلاین) و ابزارهای تستِ تعبیه‌شده در پروژه ارزیابی کنیم. این پروسه شامل سه فاز اصلی بود: اجرای تست‌های واحد، بنچمارک‌گیری هسته ریاضی، و شبیه‌سازی داده‌های \lr{CSI}.

\subsection{اجرای تست‌های واحد (\lr{Unit Tests})}
ابتدا برای اطمینان از صحت عملکرد ماژول‌های پردازشی پس از کامپایل، تست‌های نوشته‌شده در زبان \lr{Rust} را اجرا کردیم.

\begin{otherlanguage}{english}
\begin{terminalbox}
cargo test --release
\end{terminalbox}
\end{otherlanguage}

\textbf{چالش و راه‌حل:} در برخی اجراهای اولیه، به دلیل عدم دسترسی به مسیرهای پیش‌فرض دیتاست‌ها، برخی تست‌های یکپارچگی (\lr{Integration Tests}) با خطا مواجه می‌شدند. این مشکل با نادیده گرفتن تست‌های وابسته به فایل‌های خارجی و تمرکز بر تست‌های الگوریتمی (با استفاده از کتابخانه‌ی \lr{OpenBLAS} که پیش‌تر نصب کرده بودیم) با موفقیت پشت سر گذاشته شد.

\subsection{اجرای بنچمارک پردازشی (\lr{Benchmark})}
توسعه‌دهنده برای اثبات برتری سرعت \lr{Rust} نسبت به پایتون، یک فلگ اختصاصی بنچمارک در سرور قرار داده بود که ۱۰۰۰ فریم از داده‌های \lr{CSI} را به صورت فشرده پردازش می‌کرد. ما این ابزار را فراخوانی کردیم:

\begin{otherlanguage}{english}
\begin{terminalbox}
./target/release/sensing-server --benchmark
\end{terminalbox}
\end{otherlanguage}

\textbf{نتایج به دست آمده:} خروجی این دستور نشان داد که سرور توانایی پردازش علائم حیاتی و تشخیص حرکت را با سرعتی بسیار بالا (کسری از ثانیه برای ۱۰۰۰ فریم) داراست. این بنچمارک ثابت کرد که گلوگاه‌های عملکردی و کندی‌هایی که در نسخه پایتون (\lr{v1}) تجربه کرده بودیم، در این معماری به طور کامل برطرف شده‌اند.

\subsection{پردازش داده‌های \lr{CSI} از پیش ضبط‌شده (\lr{Simulation})}
برای مشاهده‌ی خروجی‌های \lr{API} و داشبورد بدون نیاز به برد \lr{ESP32}، از قابلیت شبیه‌سازی سرور استفاده کردیم. با توجه به مستندات راهنمای سرور، معماری \lr{v2} اجازه می‌دهد داده‌های منبع به حالت شبیه‌ساز (\lr{simulate}) تغییر یابند:

\begin{otherlanguage}{english}
\begin{terminalbox}
./target/release/sensing-server --source simulate
\end{terminalbox}
\end{otherlanguage}

\textbf{روند اجرا:} 
با اجرای این دستور، سرور بدون کرش کردن (و بدون جستجو برای کارت شبکه لینوکس یا برد \lr{ESP32}) بر روی پورت‌های پیش‌فرض (مانند \lr{HTTP: 8080} و \lr{WebSocket: 8765}) راه‌اندازی شد. در این حالت، سرور جریان پیوسته‌ای از داده‌های \lr{CSI} شبیه‌سازی‌‌شده را پردازش کرده و نتایج مربوط به علائم حیاتی (\lr{Vital Signs}) و نقشه حرکتی (\lr{Motion}) را بر روی \lr{API} محلی در دسترس قرار داد. این رویکرد به ما اجازه داد تا بدون درگیری با محدودیت‌های سخت‌افزاری شبکه، عملکرد الگوریتم‌های استخراج ویژگی را در نسخه \lr{Rust} به طور کامل اعتبارسنجی کنیم.

\end{document}